% Part: second-order-logic
% Chapter: syntax-and-semantics
% Section: terms-formulas

\documentclass[../../../include/open-logic-section]{subfiles}

\begin{document}

\olfileid[hi]{sol}{syn}{frm}

\olsection{पद और \printtoken{P}{formula}}

प्रथम-कोटि तर्कशास्त्र की भाँति द्वितीय-कोटि तर्कशास्त्र की अभिव्यक्तियाँ
एक मूल शब्दावली से बनती हैं, जिसमें \emph{!!{variable}},
\emph{!!{constant}}, \emph{!!{predicate}} और कभी-कभी
\emph{!!{function}} होते हैं। इनके साथ तार्किक संयोजकों, परिमाणकों और कोष्ठक
तथा अल्पविराम जैसे विराम-चिह्नों को लेकर \emph{पद} और \emph{!!{formula}}
बनाए जाते हैं। अंतर यह है कि वस्तुओं के चरों के अतिरिक्त द्वितीय-कोटि
तर्कशास्त्र में संबंधों और फलनों के चर भी होते हैं और उन पर परिमाणीकरण किया
जा सकता है। अतः द्वितीय-कोटि तर्कशास्त्र के तार्किक चिह्न प्रथम-कोटि
तर्कशास्त्र के चिह्न हैं, और इनके अतिरिक्त:

\begin{enumerate}
\item प्रत्येक $n$-स्थानिकता के द्वितीय-कोटि संबंध-!!{variable} का एक
  !!{denumerable} समुच्चय: $\Obj V_0^n$, $\Obj V_1^n$, $\Obj V_2^n$, \dots
\item द्वितीय-कोटि फलन-!!{variable} का एक !!{denumerable} समुच्चय:
  $\Obj u_0^n$, $\Obj u_1^n$, $\Obj u_2^n$, \dots
\end{enumerate}

जिस प्रकार हम प्रथम-कोटि चर $\Obj v_i$ के लिए $x$, $y$, $z$ को अधिचर के
रूप में लेते हैं, उसी प्रकार $\Obj V_i^n$ के लिए $X$, $Y$, $Z$ आदि और
~$\Obj u_i^n$ के लिए $u$, $v$ आदि को अधिचर के रूप में लेंगे।

\begin{explain}
द्वितीय-कोटि भाषा के अतार्किक चिह्न उसी प्रकार विनिर्दिष्ट किए जाते हैं जिस
प्रकार प्रथम-कोटि भाषा के: उसके !!{constant}, !!{function} और !!{predicate}
चिह्नों की सूची देकर।

प्रथम-कोटि तर्कशास्त्र में सामान्यतः !!{identity}~$\eq$ को सम्मिलित किया
जाता है। प्रथम-कोटि तर्कशास्त्र में किसी भाषा~$\Lang{L}$ के अतार्किक चिह्न
रोचक बातों को व्यक्त करने के लिए अत्यंत महत्त्वपूर्ण हैं। निस्संदेह ऐसे
!!{sentence} हैं जिनमें कोई अतार्किक चिह्न नहीं आता, किंतु केवल~$\eq$ से कोई
रोचक बात कहना कठिन है। द्वितीय-कोटि तर्कशास्त्र में संबंध-चरों और फलन-चरों
की असीमित आपूर्ति होने के कारण हम अतार्किक चिह्नों की विशेष आपूर्ति के बिना
भी वह सब कह सकते हैं जो किसी प्रथम-कोटि भाषा में कह सकते हैं।
\end{explain}

\begin{defn}[द्वितीय-कोटि पद]
~$\Lang L$ के \emph{द्वितीय-कोटि पदों} का समुच्चय $\TrmSOL[L]$,
\olref[fol][syn][frm]{defn:terms} में निम्न खंड जोड़कर परिभाषित होता है:
\begin{enumerate}
\item यदि $u$ एक $n$-स्थानी फलन-चर है और $t_1$, \dots, $t_n$ पद हैं, तो
  $\Atom{u}{t_1, \ldots, t_n}$ एक पद है।
\end{enumerate}
\end{defn}

\begin{explain}
अतः द्वितीय-कोटि पद प्रथम-कोटि पद जैसा ही दिखता है; अंतर केवल यह है कि जहाँ
प्रथम-कोटि पद में !!a{function}~$\Obj{f^n_i}$ आता है, वहाँ द्वितीय-कोटि पद
में उसके स्थान पर फलन-चर~$\Obj{u^n_i}$ आ सकता है।
\end{explain}

\begin{defn}[द्वितीय-कोटि \usetoken{s}{formula}]
भाषा~$\Lang L$ के \emph{द्वितीय-कोटि !!{formula}} का समुच्चय
~$\FrmSOL[L]$, \olref[fol][syn][frm]{defn:formulas} में निम्न खंड जोड़कर
परिभाषित होता है:
\begin{enumerate}
\item यदि $X$ एक $n$-स्थानी विधेय-चर है और $t_1$, \dots, $t_n$
  भाषा~$\Lang L$ के द्वितीय-कोटि पद हैं, तो
  $\Atom{X}{t_1,\ldots, t_n}$ एक परमाणु !!{formula} है।

\tagitem{prvAll}{यदि $!A$ !!a{formula} है और $u$ फलन-चर है, तो
  $\lforall[u][!A]$ !!a{formula} है।}{}

\tagitem{prvAll}{यदि $!A$ !!a{formula} है और $X$ विधेय-चर है, तो
  $\lforall[X][!A]$ !!a{formula} है।}{}

\tagitem{prvEx}{यदि $!A$ !!a{formula} है और $u$ फलन-चर है, तो
  $\lexists[u][!A]$ !!a{formula} है।}{}

\tagitem{prvEx}{यदि $!A$ !!a{formula} है और $X$ विधेय-चर है, तो
  $\lexists[X][!A]$ !!a{formula} है।}{}
\end{enumerate}
\end{defn}

\end{document}
